% =============================================================================
%  EXPANDED APPENDIX  --  \input after \appendix in main.tex
%  No preamble, no \begin{document}. Sequence of \section blocks only.
%  Notation matches main.tex (z, A, S_g, u, D, Eq.~\ref{eq:doop}).
% =============================================================================

\section{Model and objective in full}\label{app:model_full}
This section gives the world model and its training objective at the level of exact
shapes and hyperparameters. All shapes below are those of the deployed world model at
epoch $60$. The gene axis has $G{=}4366$ genes and the latent has $d{=}357$ gene
programs. The decoder weight has shape $4366\times357$, the latent directed graph $A$ has
shape $357\times357$, and the per-TF label matrix has shape $4366\times100$.

\subsection{Encoder and latent variables}
The encoder is an amortized Gaussian. It reads a cell expression vector
$x\in\mathbb{R}^{4366}$ concatenated with a $16$-dimensional modality embedding
(Joung is modality $0$, Briggs is modality $1$) and maps the $4382$-dimensional input
through a hidden layer of width $512$ to a mean $\mu$ and log-variance
$\log\sigma^2$ over the exogenous term $\epsilon\in\mathbb{R}^{357}$. The first encoder
weight therefore has shape $512\times4382$ and the modality embedding table has shape
$2\times16$. We sample $\epsilon$ with the standard reparameterization. The
observational forward pass solves the latent structural equation model (SEM) for $z$
and decodes $x{=}Dz{+}b$.

\subsection{The latent SEM and its closed-form solve}
We model the resting latent state by a linear SEM over a learned signed directed graph
$A\in\mathbb{R}^{357\times357}$ driven by $\epsilon$,
\begin{equation}
z \;=\; A^{\top} z + \epsilon
\;\;\Longleftrightarrow\;\;
(I-A^{\top})\,z=\epsilon
\;\;\Longrightarrow\;\;
z=(I-A^{\top})^{-1}\epsilon .
\label{eq:sem_solve}
\end{equation}
The closed form follows because $I-A^{\top}$ is invertible whenever the spectral
radius of $A$ is below $1$. We enforce that by a power-iteration spectral-norm clamp
that rescales every sampled $A$ so $\lVert A\rVert_2 \le 0.85$, which guarantees
$I-A^{\top}$ is nonsingular and every downstream solve is well posed. The solve itself
is a direct linear solve with a sanitized escalating-ridge fallback for the rare
early-training singular case. The adjacency is a factored Gumbel-sigmoid gate times a
real-valued weight times a no-self-loop mask,
$A_{ij}=\mathrm{gate}_{ij}\cdot w_{ij}\cdot(1-\delta_{ij})$, with
$\mathrm{gate}_{ij}=\sigma((\theta_{ij}+\text{Gumbel})/\tau)$ in training and the hard
threshold $\mathbf{1}[\sigma(\theta_{ij}/\tau)>0.5]$ at evaluation. The temperature
$\tau$ anneals from $1.0$ to $0.25$ over training. At evaluation the graph is therefore
deterministic and signed.

\subsection{The shift-mutilate do-operator and $k$-TF composition}
A perturbation of a transcription factor (TF) $g$ enters as a soft intervention with two
parts. First it adds a program-space shift $S_g\in\mathbb{R}^{357}$. Second it applies a
per-program mutilation gate $u_g\in[0,1]^{357}$ that partially severs incoming edges. For
a set $\mathcal{P}$ of perturbed TFs the deployed operator sums the single-TF shifts,
injects the sum into the exogenous term, forms one aggregate gate, and does one mutilated
solve. This is Eq.~\eqref{eq:doop} of the main text. Concretely, in the structural
intervention mode used throughout,
\begin{align}
S &= \textstyle\sum_{g\in\mathcal{P}} S_g \;(+\,S_{\mathrm{int}}), &
\tilde{\epsilon} &= \epsilon_0 + S, \label{eq:doop_struct}\\
u &= \operatorname*{mean}_{g\in\mathcal{P}} \operatorname{clip}\!\big(\beta\lvert S_g\rvert,\,0,\,1\big),\ \beta{=}1, &
z_{\mathrm{post}} &= \big(I-\operatorname{diag}(1-u)A^{\top}\big)^{-1}\tilde{\epsilon}. \nonumber
\end{align}
Two design points matter. The gate $u$ is the per-single mean of $\operatorname{clip}(\beta\lvert S_g\rvert)$
rather than $\operatorname{clip}(\beta\lvert\sum_g S_g\rvert)$. The summed form saturates
about $41\%$ of programs to $u{=}1$ on combinations, which detaches $A$ and degrades
combination prediction, so we use the per-single mean, which is identical for a single TF.
In structural mode the shift is injected only into
$\tilde{\epsilon}$, so any non-target program moves only through $A$. If $A{=}0$ there is
no downstream effect, which is what makes $A$ load-bearing for prediction by construction.
The operator reduces to the additive shift $z_{\mathrm{prev}}+z_S$
when $u\to1$ on all driven programs and to predict-control when $u\to0$. Non-additivity
for combinations arises because $u$, and hence $M=\operatorname{diag}(1-u)$, depends
nonlinearly on the perturbed set. An optional symmetric low-rank interaction head adds a
combination-specific shift $S_{\mathrm{int}}=W_{\mathrm{out}}\big(\tfrac12((\sum_i x_i)^2-\sum_i x_i^2)\big)$
with $x_i=U(\text{ESM}_i)$ and an interaction rank of $32$. $S_{\mathrm{int}}$ is $0$ for a
single gene. The operator is validated against a naive fixed-point solve. A $k$-TF
prediction is therefore one linear solve assembled from
the same single-TF shifts used for singles, exact and zero-shot for novel combinations.

\subsection{Per-TF shift from ESM embeddings}
The shift $S_g$ is produced by a learned shift encoder from a per-gene descriptor
$\ell_g$. The descriptor is an ESM2 (Evolutionary Scale Modeling v2) protein language
model embedding of the TF, reduced by principal component analysis and z-scored to $100$
dimensions. The descriptor exists for
any TF in the panel, including held-out TFs, which is what lets the shift generalize
zero-shot. The encoder is a one-hidden-layer multilayer perceptron (MLP) of width $1024$
with a LeakyReLU$(0.2)$ nonlinearity mapping $100\to357$, with a raw real-valued output and
no output activation. Its first weight has shape $1024\times100$ and its output weight has
shape $357\times1024$. A second head shares the trunk and emits per-program
gate logits, so the shift is $S_g=c\cdot(\text{gate}\odot\text{raw})$ with a learned scalar
penetrance $c$ and a sparse target gate that is Gumbel-sigmoid in
training and hard at evaluation. The gate open
probability is penalized by an $L_0$ target-sparsity term. A leakage-safe train-only
retrieval memory over embedding keys can add an embedding-similarity-weighted shift with
the query row self-masked, which we keep off by default in the reported model.

\subsection{The masked gene-program (SENA) decoder}
The decoder is linear, $x{=}Dz{+}b$ with $D\in\mathbb{R}^{4366\times357}$. Each latent
column is tied to a Gene Ontology Biological Process (GO:BP) program by a membership mask
$M\in\mathbb{R}^{4366\times357}$. We refer to this masked gene-program decoder scheme as
SENA, following \citep{sena2025}. The effective decoder weight is the masked-learnable
product $D\odot M$, so annotated gene-program pairs learn at full strength while
off-annotation pairs are shrunk by a leak factor $\lambda$. This forces gene covariance to
route through the program latents and $A$ rather than through arbitrary gene-gene edges.
The mask is built from the Enrichr GO Biological Process 2023 library. Candidate terms are
GO:BP sets with at most $30$ genes and at least $5$ genes present on our axis. Terms are
greedily de-duplicated largest-first so no kept term overlaps a previously kept one by more
than $0.5$ under the containment overlap
$\lvert A\cap B\rvert/\min(\lvert A\rvert,\lvert B\rvert)$. Off-annotation entries take the
leak value $\lambda{=}0.3$. The GO:BP selection yields $357$ kept programs on our axis, so
the latent dimension is set to the number of kept programs and the deployed model has
$d{=}357$ program latents. The decoder mask is not frozen, so $D$ is masked-learnable
rather than fixed to the mask.

\subsection{Training objective, prior-off}
The total loss is
\begin{equation}
\begin{split}
\mathcal{L}
={}& \lambda_{\mathrm{rec}}\,\mathcal{L}_{\mathrm{rec}}
+ \beta_{\mathrm{eff}}\,\mathcal{L}_{\mathrm{KL}}
+ \lambda_{\mathrm{pert}}\,\mathcal{L}_{\mathrm{pert}}
+ \lambda_{\mathrm{spec}}\,\mathcal{L}_{\mathrm{spec}}\\
&+ \lambda_{\mathrm{ortho}}\,\mathcal{L}_{\mathrm{ortho}}
+ \lambda_{\mathrm{sparse}}\,\mathcal{L}_{\mathrm{sparse}}
+ \lambda_{\mathrm{prior}}\,\mathcal{L}_{\mathrm{prior}} .
\end{split}
\label{eq:total_loss}
\end{equation}
We train prior-off. The graph prior, decoder-orthogonality, graph-sparsity, decoder-$L_1$,
and acyclicity weights are all zero
($\lambda_{\mathrm{prior}}{=}0$, $\lambda_{\mathrm{ortho}}{=}0$, $\lambda_{\mathrm{sparse}}{=}0$,
$\lambda_{\mathrm{D\text{-}sparse}}{=}0$, and both acyclicity coefficients $0$),
so no external transcription-factor--target database edge is ever used to shape $A$. The
acyclicity term is skipped entirely when its weights are zero to avoid $0\times\text{NaN}$
poisoning. The training log confirms the acyclicity residual stays at $0$ throughout.
Table~\ref{tab:loss_terms} lists every active term with its exact form and weight for the
deployed model.

The reconstruction term is a plain mean squared error (MSE), modality-gene-masked so a
cell only scores genes its modality observes. The Kullback--Leibler (KL) divergence is the
standard diagonal-Gaussian variational term with a free-bits floor, and its weight $\beta$
ramps linearly from $0.01$ to $0.5$ over the first $30$ epochs. The perturbation term
$\mathcal{L}_{\mathrm{pert}}$ (weight $\lambda_{\mathrm{pert}}{=}10$, clamped at $50$ before
scaling) is the sum of the sub-terms in Table~\ref{tab:loss_terms}, computed on cell
sentences of $S{=}64$ cells so the mean-shift target is a per-sentence pseudobulk $\Delta$.
The DE-gene MSE sub-term deliberately preserves magnitude and is not whitened, and is
restricted to observed differentially expressed (DE) genes above the threshold $0.1$, which
is the term that carries per-gene response magnitude.

\begin{table}[h]\centering\footnotesize
\caption{Active training-loss terms for the deployed world model. All weights are exact
configured values. Terms with weight $0$ are omitted (graph prior, decoder-orthogonality,
graph-sparsity, decoder-$L_1$, acyclicity, energy optimal-transport, and shift-$L_1$).
MMD denotes maximum mean discrepancy and InfoNCE the info-noise-contrastive estimation
loss.}
\label{tab:loss_terms}
\begin{tabular}{@{}llll@{}}
\toprule
Term & Form (drift $=$ mean train-single $\Delta$) & Weight & Value \\
\midrule
Reconstruction & $\lVert (\hat{x}-x)\odot m_{\mathrm{mod}}\rVert^2$ mean-MSE & $\lambda_{\mathrm{rec}}$ & $1.0$ \\
KL & $-\tfrac12\overline{(1+\log\sigma^2-\mu^2-\sigma^2)}$, free-bits & $\beta$ & $0.01\!\to\!0.5$ \\
Perturbation MSE & $\lVert \Delta_{\mathrm{pred}}-\Delta_{\mathrm{obs}}\rVert^2$ & $\lambda_{\mathrm{MSE}}$ & $1.0$ \\
Distribution matching (MMD) & Gaussian-RBF MMD$^2$, median bandwidth, per sentence & $\lambda_{\mathrm{MMD}}$ & $1.0$ \\
Goal specificity (InfoNCE) & symmetric NCE on drift-removed norm.\ $\Delta$, temp $0.2$ & $\lambda_{\mathrm{NCE}}$ & $1.0$ \\
DE-gene MSE & unwhitened MSE on $\{|\Delta_{\mathrm{obs}}{-}\mathrm{drift}|{>}0.1\}$ genes & $\lambda_{\mathrm{DEG}}$ & $8.0$ \\
Target sparsity ($L_0$) & $\overline{\sum_k \text{gate-prob}_k}$, ramped & $\lambda_{L_0}$ & $0.03$ \\
Spectral & $\operatorname{relu}(\sigma_{\max}(A)-0.95)$ & $\lambda_{\mathrm{spec}}$ & $0.1$ \\
\bottomrule
\end{tabular}
\end{table}

A gradient-reversal dataset classifier on the encoder mean $\mu$ drives latent
modality-invariance and is added with an implicit coefficient of $1.0$ when there is more
than one modality. The mutilation is run in shift mode with $\beta{=}1.0$. The mutilation
clip of $0.5$ is inert in shift mode, where the clamp is to $1.0$.

% =============================================================================
\section{Datasets, splits, and checkpoint discipline}\label{app:data}
\subsection{Sources}
We use three single-cell resources. The Joung TF atlas is a human embryonic stem cell
(hESC) single-TF overexpression atlas at day $7$ \citep{joung2023tfatlas}. It contributes
single-TF perturbations plus a small number of arrayed combinations. The Briggs
hESC$\to$fate overexpression screen is a dense combinatorial TF screen read out at days
$14$ and $28$, deposited at GEO under accession GSE328457 \citep{briggs2026}. It
contributes high-order TF combinations and pluripotent hESC-start controls. The Cao fetal
cell atlas defines the target-fate marker signatures \citep{cao2020fetal}.

\subsection{Gene axis and modality alignment}
The shared gene axis is the Joung $4366$-gene vocabulary. Briggs cells are aligned onto it
by matching counts-per-$10{,}000$ log-normalized values on shared genes and zero-filling
the rest. About $2800$ of the $4366$ genes are observed in Briggs. A per-modality gene mask
keeps the vocabulary at $4366$ but scores the reconstruction loss only on genes each
modality observes. Joung sees all $4366$ genes, Briggs only its observed subset. Each
Briggs cell carries an in-vocabulary TF set obtained by collapsing isoforms and keeping the
modelable TFs. Cells are grouped by their in-vocabulary TF set. Sets of in-vocabulary size
$2$ to $8$ with at least $5$ replicate cells become combination
perturbations. The per-cell median in-vocabulary set size is about $5$ and reaches into the
teens, so Briggs supplies the high-order regime.

\subsection{Answer-key holdout and leakage-free splits}
The answer key is held out of training so recovery is honest. A Briggs combination is held
out if it shares at least two members with $P_1$ or $P_2$, is one of the six canonical hematopoietic stem and progenitor cell (HSPC)
pairs, or is one of the arrayed-validated combinations from the supplementary table. The
two validated cocktails are
$P_1=\{\text{TAL1, RPA3, ZFP90, RNF138, LMO2, KLF2}\}$ and
$P_2=\{\text{DDIT3, MZF1, HOXB3, MIER1, SIRT6, ZNF85}\}$. In the deployed
model this holds $222$ answer-key combinations out of training while keeping them for
evaluation. Single TFs are
split with combo-involved TFs forced into training, so a single-TF holdout never silently
guts combination coverage.
Fate targets are split by fate name into folds to prevent cross-source leakage.

\subsection{The two deployed checkpoints}
Two checkpoints back the paper, both trained for $60$ epochs.

\textbf{Deployed production world model.} Trained with both the single-TF and combination
training fractions set to $1.0$, so every single TF and every non-answer-key combination is
in training and only the $222$ answer-key combinations are held out. Like the held-out
sibling, it warm-starts from a singles-only base model and oversamples single-TF sentences
$8\times$ during the combination finetune, so single-TF effects are retained. The training
log records $963{,}974$ cells, a per-epoch reconstruction subsample of $150{,}000$ cells,
$495$ single-TF blocks and $1622$ combination blocks ($2117$ blocks total), and $22{,}328$
cell sentences per epoch. This is the checkpoint used for all causal-structure validation
(App.~\ref{app:model_full} numbers and Sec.~\ref{sec:causal}) and as the frozen reward
oracle for the planner. Because its training fractions are $1.0$, a held-out-singles score
on this checkpoint measures memorization and is not a genuine held-out number. Genuine
held-out single and combination prediction uses the next checkpoint.

\textbf{Held-out combination finetune.} Trained with a single-TF training fraction of
$0.85$ and a combination training fraction of $0.8$ under a fixed split seed, warm-started
from a singles-only base model, with recency-weighted replay of combinations during the
finetune. This split yields $432$ train and $63$ held-out Joung singles, $36$ train and $8$
held-out Joung combinations, and Briggs combinations split into train and held-out plus the
$222$ answer-key combinations. Two additional seed variants provide the three seeds behind
the main-text prediction $p$-values. The main-text held-out advantages of $+0.184$
(answer-key combinations) and $+0.225$ (held-out combinations) in drift-removed
\PearsonDEG are computed on this checkpoint across the three seeds. The epistasis
panels (Fig.~\ref{fig:epi}A--D) also use this held-out sibling for the combination-trained
model and its singles-only base for the singles-only arm, so the interaction claims are
evaluated on the same recipe as the planner world model. Only the enrichment panel
(Fig.~\ref{fig:epi}E) uses the deployed production world model.

\subsection{Fate-target annotation}\label{app:fate}
Target fates are defined by a marker-signature pipeline over the Cao fetal cell atlas. We
start from the Cao cell-type mean log-expression profiles, restrict to genes on our shared
axis, and center them. Cell types are merged into fate groups by average-linkage
hierarchical clustering at cosine similarity $0.78$ on the centered profiles, and each fate
group is named by its dominant tissue--type suffix. For each fate group we take the top
$50$ genes by one-vs-rest over-expression as its marker set and form a normalized marker
signature.

Every candidate perturbation (Briggs combinations, Joung single-TF overexpressions, and
Joung arrayed combinations) is summarized by its drift-removed mean response, the residual
centroid $\Delta$ obtained by subtracting the source-level mean $\Delta$. Each perturbation's
residual $\Delta$ is scored against all marker signatures and assigned to its argmax fate. A
fate is \emph{kept} only if it is \emph{target-separable}, meaning its marker signature is
isolated from every other signature (max cosine similarity below $0.75$), and
\emph{reachable}, meaning it collects at least a minimum number of marker-argmax member
cells ($5$ for Briggs, $3$ for Joung). Each kept fate's target is the drift-removed mean
$\Delta$ (residual centroid) of its marker-argmax member cells. Kept fates are then tiered.
Tier~A fates additionally have a positive per-fate cosine silhouette in the marker-gene
(specific-gene) subspace, and Tier~B fates have a non-positive silhouette. Crucially, the
annotation does \emph{not} gate on silhouette. Both tiers are retained, and the silhouette
only labels the clean-answer-key core (Tier~A) versus the remainder (Tier~B). Reachable-fate
confidence uses a proper Benjamini--Hochberg correction, which moves the fate set by only one
or two Briggs fates relative to a raw permutation cutoff.

Each fate additionally stores its Cao marker genes, its member count, its driver TFs, and its
answer-key combinations. The driver and answer-key lists derive from the Briggs $P_1$/$P_2$
canonical cocktails and the arrayed-combo table (Sec.~\ref{app:data}). The pipeline yields
$84$ fates, namely $65$ Briggs, $16$ Joung single-TF, and $3$ Joung combination fates
(Table~\ref{tab:fates}). The zero-shot HSPC fate,
the Briggs liver hematopoietic-stem-cell fate (Tier~B), is forced to fold $-1$ and is always
held out. Its evaluation uses the dedicated Briggs $P_1\cup P_2$ answer key rather than Cao
markers.

Fates are split into a five-fold leave-fates-out cross-validation partition by \emph{fate
name}, so all sources of one fate name share a fold and there is no cross-source leakage. The
folds are stratified by name-level tier and assigned round-robin by descending member count,
so each Tier~A fate is held out exactly once. In each cross-validation fold $f\in\{0,\dots,4\}$
the fates assigned to fold $f$ are held out for evaluation and the rest are used for
training, and the single fate at fold $-1$ (HSPC) is held out in every fold. The per-fold
held-out counts are $19$, $18$, $14$, $16$, $16$ for folds $0$--$4$ plus the $1$ always-held-out
HSPC fate at fold $-1$.

\begin{table}[h]\centering\footnotesize
\caption{Fate-target counts by source and tier. Tier~A fates additionally clear a positive
marker-subspace silhouette (clean answer-key core). Tier~B fates are target-separable and
reachable but silhouette-negative. The HSPC fate (the Briggs liver hematopoietic-stem-cell
fate) is one of the $65$ Briggs fates, is Tier~B, and is forced to fold $-1$ (always held
out).}
\label{tab:fates}
\begin{tabular}{@{}lrrr@{}}
\toprule
Source & Fates & Tier A & Tier B \\
\midrule
Briggs (combinations)      & $65$ & $10$ & $55$ \\
Joung single-TF            & $16$ & $2$  & $14$ \\
Joung combination          & $3$  & $2$  & $1$  \\
\midrule
Total                      & $84$ & $14$ & $70$ \\
\bottomrule
\end{tabular}
\end{table}

% =============================================================================
\section{Training and planner hyperparameters}\label{app:hparams}
\subsection{World-model hyperparameters}
Table~\ref{tab:wm_hparams} lists the exact world-model settings for the deployed model.
Training runs in three phases inside the $60$ epochs, recorded in the training log. Phase A
is epochs $1$ to $10$, reconstruction only with the directed graph frozen. Phase B is epochs
$11$ to $50$, directed-graph learning active with the full loss. Phase C is epochs $51$ to
$60$, full loss with temperature annealing. The world model trains on CPU high-memory nodes
with about $340$ GB of RAM.

The architecture is the model of App.~\ref{app:model_full}: a per-cell encoder, a latent
SEM over $d{=}357$ GO:BP program latents solved in closed form, an ESM2 per-TF shift encoder
($100$-dim descriptor $\to$ $1024$ hidden $\to$ $357$), the shift-mutilate do-operator for
$k$-TF composition, and the masked-GO (SENA) linear decoder to the $4366$-gene axis. These
same settings train both deployed checkpoints described in Sec.~\ref{app:data} (``The two
deployed checkpoints''). The deployed production world model is the epoch-$60$ checkpoint,
used for all causal-structure validation and as the frozen planner reward oracle. The
held-out combination finetune shares every model hyperparameter in
Table~\ref{tab:wm_hparams} and differs only in its train fractions (single-TF fraction
$0.85$, combination fraction $0.8$) and its singles-only-base warm-start. It backs the
held-out prediction and epistasis panels. Both checkpoints warm-start from a singles-only
base and oversample single-TF sentences $8\times$ during the combination finetune.

\begin{table}[h]\centering\footnotesize
\caption{World-model hyperparameters for the deployed model. Values not present in the
configuration are marked not specified.}
\label{tab:wm_hparams}
\begin{tabular}{@{}llll@{}}
\toprule
Optimizer & Adam & Gene axis $G$ & $4366$ \\
Learning rate & $10^{-3}$ & Latent programs $d$ & $357$ (GO:BP mask) \\
LR warmup epochs & $5$ & Decoder & masked-GO (SENA), $\lambda{=}0.3$ \\
Epochs & $60$ & Shift descriptor & ESM2, $100$-dim \\
Batch size & $256$ & Shift MLP & $100\!\to\!1024\!\to\!357$, LeakyReLU$(0.2)$ \\
Grad accumulation & $4$ & Interaction rank & $32$ \\
KL $\beta$ & $0.01\!\to\!0.5$ over $30$ ep & Graph samples (train/eval) & $1$ / $32$ \\
$\lambda_{\mathrm{pert}}$ & $10$ & Mutilation $u_{\max}$, $\beta$ & $0.5$, $1.0$ \\
DEG-MSE weight, thr & $8.0$, $0.1$ & Spectral weight & $0.1$ \\
Contrastive temp & $0.2$ & Prior / sparse / ortho / acyclicity & $0$ (prior-off) \\
Cell-sentence $S$ & $64$ & Sentences per block & $4$ \\
Sentence batch & $32$ & Single-sentence mult & $8$ \\
Recon subsample & $150{,}000$/epoch & $\tau$ anneal & $1.0\!\to\!0.25$ \\
\bottomrule
\end{tabular}
\end{table}

\subsection{The GFlowNet planner}
The deployed planner is a trajectory-balance generative flow network (GFlowNet) over the frozen world model. A cocktail
is built one TF at a time by an autoregressive forward policy. Because the construction is
autoregressive there is a single parent per state, so the backward policy is trivially
$P_B\equiv1$ and the trajectory-balance loss for a terminal cocktail $x$ under goal $g$ is
\begin{equation}
\mathcal{L}_{\mathrm{TB}}
=\Big(\log Z(g) + \textstyle\sum_t \log P_F(a_t\mid s_t,g) - \beta\cdot\Phi(x,g)\Big)^2 ,
\label{eq:tb}
\end{equation}
with reward $R(x)=\exp(\beta\,\Phi(x,g))$ \citep{bengio2021gflownet} and reward exponent
$\beta{=}48$. The deployed flagship uses a \emph{goal-conditioned} partition function
$\log Z(g)$, where a per-goal head absorbs the goal-dependent normalizer so the forward policy is
free to spread mass across the many high-reward cocktails for each goal. This is the single
most important lever for diversity. The alternative of a single shared scalar $\log Z$ for
all goals forces the forward policy to carry the goal-dependent normalizer, and it collapses
to an amortized argmax at convergence (effective number of distinct sampled sets
$\mathrm{ess}_q\approx1$--$2$, versus up to the target ceiling of $\approx46$ under the
conditional head). We report the shared-$\log Z$ variant as an ablation.

\paragraph{Ensemble-pessimistic reward.} The potential $\Phi(x,g)$ is the drift-removed match
of the world-model-predicted state to the target fate. To guard against exploiting any single
world model's errors, the reward is \emph{ensemble-pessimistic}. We evaluate the same cocktail
under $K{=}6$ independently trained world models and combine their potentials as
\begin{equation}
\Phi_{\mathrm{pess}}(x,g)=\operatorname*{mean}_{k=1}^{K}\Phi^{(k)}(x,g)\;-\;\lambda\cdot\operatorname*{std}_{k=1}^{K}\Phi^{(k)}(x,g),
\qquad \lambda=0.25 ,
\label{eq:pess}
\end{equation}
so cocktails on which the ensemble disagrees are penalized. The pessimism weight $\lambda$
trades recovery against sampling diversity. The deployed value is $\lambda{=}0.25$.

Off-policy replay mixes a fraction $\gamma{=}0.5$ of each batch from a recency-weighted buffer
of past cocktails (capacity $5000$), which is valid for trajectory balance. The reward oracle
is the frozen world-model ensemble, so no surrogate or acquisition proxy is needed. Batches
are $16$ trajectories per update, trained for $10^{6}$ iterations. The per-source action
budget is set at reset, so Joung caps at $k{\le}2$ and Briggs at $k{\le}8$. Prioritized,
reward-weighted replay is a reported negative (App.~\ref{app:rl}). The deployed setting is
plain recency-weighted replay.

The policy is a set-based autoregressive policy over the vocabulary
$\{\text{STOP}\}\cup\text{candidates}$. A shared encoder feeds all slots, so cost grows
linearly in $k$ not combinatorially. STOP is forbidden at slot $0$ and already-chosen genes
are masked at later slots. The action head is a plain multilayer perceptron (MLP)
over the per-slot features.

\subsection{Design selection}
The deployed flagship, the goal-conditioned GFlowNet with the plain MLP head and
$\lambda{=}0.25$, was selected by a sweep over the partition function (shared versus
goal-conditioned $\log Z$) and the ensemble-pessimism weight
$\lambda\in\{0.1,0.25,0.5,0.75,1.0\}$, at $10^6$ iterations across five folds and three seeds
on the frozen production world model with a curated candidate pool and per-source $k$ caps.
The goal-conditioned $\log Z$ is what enables diverse sampling and $\lambda{=}0.25$ gives the
best recovery--diversity trade-off (App.~\ref{app:rl}). The InfoNCE fate-specificity term uses
scale $\beta_{\mathrm{contrast}}{=}0.3$; a no-contrast control that drops it is reported in
App.~\ref{app:rl}.

% =============================================================================
\section{Evaluation metrics}\label{app:metrics}
This section defines every metric precisely. For a planner the per-fate output is a
frequency dictionary of proposed TFs over the best-of-$N$ samples. Picks are ranked by
descending selection frequency.

\textbf{Driver recall at $k$.} With driver set (answer key) $\mathcal{D}$ and the top-$k$
ranked picks $\mathcal{T}_k$, $\mathrm{drv\_recall}@k = \lvert\mathcal{T}_k\cap\mathcal{D}\rvert/\lvert\mathcal{D}\rvert$.
The denominator is the answer-key size, so this is recall not precision. We report
$k\in\{5,10,20\}$.

\textbf{Precision at $k$.} $\mathrm{prec}@k = \lvert\mathcal{T}_k\cap\mathcal{D}\rvert/\min(k,\lvert\mathcal{T}_k\rvert)$.
We report $k\in\{5,10\}$.

\textbf{Cross-fate discrimination, own minus other.} For a scorer $s$, discrimination is the
score of a fate's picks against its own answer key minus the mean score of the same picks
against every other held-out fate's answer key, $\mathrm{disc}=s(\text{own})-\overline{s(\text{other})}$.
We report $\mathrm{disc\_dr20}$ using recall@$20$ and
$\mathrm{disc\_soft1}$ using SOFT-1. The oracle plays the true drivers as picks. For
$\mathrm{disc\_dr20}$ the oracle equals $1-\overline{\text{cross-fate driver-set overlap}}$,
so a high oracle means the answer-key driver sets are nearly disjoint across
fates and a perfect selector separates them, while a low oracle means fates share drivers
and cap what any planner can achieve. A random planner sits near zero because its picks hit
own and other drivers equally.

\textbf{SOFT-1 and SOFT-2.} These score in a functional space built from drift-removed,
unit-normed single-TF $\Delta$ vectors with $k$-means module labels. SOFT-1 is the
frequency-weighted mean, over picks, of
the maximum cosine of a pick's single-TF $\Delta$ to any driver's single-TF $\Delta$. SOFT-2
is the frequency-weighted fraction of picks whose $k$-means module coincides with any driver's
module. The variants m12, m24, m48 use $12$, $24$, or $48$ modules. These metrics saturate.
The nearest-cosine over a broad answer set in a low-dimensional unit-normed space means even
random TFs land close to some driver and share one of only twelve modules, so a matched-random
null already scores about $0.28$ to $0.36$ and the policy barely separates. For this reason
SOFT-1 and SOFT-2 are reported in the appendix with their null values shown, and the
headline uses driver recall, precision, and the own-minus-other discrimination forms.

\textbf{Best residual match.} For each sampled cocktail the world model applies the
do-operator, decodes, subtracts the control reconstruction and the drift, and matches the
drift-removed predicted $\Delta$ against the target on the active dimensions. The best
residual match is the maximum match over the best-of-$N$ sampled cocktails.

\textbf{The matched-random null.} For each held-out fate the null records the policy's
per-sample cocktail sizes, then draws for each sample a random cocktail of the same size
from the candidate pool and recomputes every metric. The reported null value is the mean
over $R$ permutations ($R{=}500$ for the reported runs). When a policy is gated to a
per-goal shortlist, the reported evaluation draws the null from the full $280$-candidate
pool, which gives a pipeline-versus-ungated contrast consistent across all sources.

\textbf{The $1000$-shuffle-$A$ permutation test.} For the causal-structure claim we ablate
$A$ inside the frozen do-operator. The variants are original, reversed ($A^{\top}$),
symmetrized ($\tfrac12(A+A^{\top})$), and shuffled (a permutation of the off-diagonal
entries, which preserves the value multiset and edge count). Every non-original variant is
rescaled to
match $\lVert A\rVert_2$, so magnitude is not a confound. The null draws $N{=}1000$ shuffled
$A$ matrices, each with a fresh seed, and stores the mean drift-removed \PearsonDEG
across perturbations. The single-TF prediction metric is the per-perturbation Pearson of the
drift-removed predicted residual against the drift-removed observed residual on the set of
DE genes above the threshold $0.05$, with the perturbed TF itself excluded so the metric
measures downstream recovery. The permutation $p$-value is the add-one estimate
$(1+\#\{\text{null}\ge\text{original}\})/(N+1)$, which gives $1/1001\approx9.99\times10^{-4}$
when no shuffle exceeds the original.

\textbf{Paired Wilcoxon across fates and Benjamini--Hochberg.} The primary planner
significance test is a one-sided paired Wilcoxon signed-rank of the policy against its
matched-random null across held-out fates. The replicate is the fate when at least two
held-out fates exist, with per-fate values
averaged over seeds. For a single-fate source such as HSPC the test falls back to pairing
across the fold-by-seed training replicates, which gives up to $15$ paired points for five
folds and three seeds. Benjamini--Hochberg correction is applied across the whole
source-by-metric grid. The same one-sided paired Wilcoxon signed-rank test (dropping zero
differences) is used for the $66$-TF and $14$-combination causal-structure tests.

\textbf{Standard error of the mean (with one delta-degree-of-freedom).} Standard errors of the
mean, abbreviated s.e.m. in the tables and distinct from the structural equation model above,
use the per-seed means as the replicate axis. Within each seed, per-fate values are pooled
across folds and averaged, and the standard error of the mean is the sample standard
deviation of those per-seed means with one delta-degree-of-freedom
divided by $\sqrt{n_{\mathrm{seeds}}}$. A single replicate is reported as undefined rather
than zero.

\textbf{Baselines.} MPC-exact is a model-predictive control (MPC) world-model beam search (beam width $8$) that
forward-selects the cocktail maximizing the predicted drift-removed match. Additive-greedy
is a forward selection that sums observed drift-removed single-TF $\Delta$ vectors with no
world-model interaction. Random-shortlist is a single
random pick from the goal shortlist. Every baseline is restricted to the same per-goal
shortlist as the policy, so the comparison isolates selection within the pool.

% =============================================================================
\section{Full cross-arm planner evaluation}\label{app:rl}
The deployed planner is the goal-conditioned GFlowNet with the plain MLP head and the
ensemble-pessimistic reward at $\lambda{=}0.25$ (Eq.~\eqref{eq:pess}), trained over three
folds and five seeds. We evaluate it on two axes against non-learned baselines, asking (1) whether it
\emph{recovers} the true driver TFs of a held-out fate, and (2) whether it \emph{samples
diverse, reward-proportional} cocktails, with a zero-shot HSPC case study. The genuine wins
are that the planner is amortized and cheap, is diverse-at-quality, and on the zero-shot HSPC fate
covers more of the arrayed answer key with biologically coherent regulators than search does. On
per-driver recovery it reaches \emph{parity} with model-predictive beam search, not a win.

\paragraph{Why a GFlowNet trainer.} Against a one-step set-PPO trainer, the GFlowNet samples
many distinct high-quality cocktails while set-PPO collapses toward a single mode, which is
why the deployed planner is a GFlowNet rather than a policy-gradient trainer.

\subsection{Axis 1, driver recovery}
On the finalized flagship we report held-out Briggs recovery under three readouts of the
$256$ sampled cocktails, each against its \emph{matched} null (Table~\ref{tab:flagship_rec}).
The \emph{raw} readout ranks TFs by sampling frequency and is compared to a matched-breadth
permutation null ($\approx0.072$) that draws random size-matched cocktails and ranks them the
same way. The \emph{elite-$10$} readout keeps the top $10\%$ of the $256$ by reward and ranks
within that set, and the \emph{leave-one-out} (LOO) readout ranks each TF by its marginal
reward contribution. Both are compared to the \emph{random-plus-elite} null ($\approx0.42$),
which passes random cocktails through the same reward filter and so measures what the reward
filter alone buys. Using the weak permutation null on an elite readout would over-credit the
policy for the filter's work, so each readout is matched to its own null. Model-predictive
beam search is likewise a reward-guided selector, so the fair comparison between the policy
and beam search uses the common random-plus-elite null.

\begin{table}[h]\centering\footnotesize
\caption{Deployed flagship (goal-conditioned GFlowNet, MLP head, ensemble pessimism
$\lambda{=}0.25$), held-out Briggs driver recall at $20$, three folds $\times$ five seeds.
Each readout is scored against its matched null. Here $\Delta$ is the policy minus that null.
World-model (WM) evaluations per fate index inference cost. The random-plus-elite null is the
common reward-guided baseline used to compare the policy against beam search.}
\label{tab:flagship_rec}
\begin{tabular}{@{}llccc@{}}
\toprule
Planner / readout & null & abs r@$20$ & $\Delta$ vs null & WM evals/fate \\
\midrule
policy, raw            & permutation ($0.072$)   & $0.455$ & $+0.382$ & $256$ \\
policy, elite-$10$     & random$+$elite ($0.38$) & $0.582$ & $+0.200$ & $256$ \\
policy, LOO            & random$+$elite ($0.42$) & $0.712$ & $+0.291$ & $207$ \\
MPC-exact (beam $8$)   & random$+$elite ($0.39$) & $0.507$ & $+0.120$ & $9240$ \\
additive               & permutation ($0.072$)   & $0.012$ & $-0.060$ & --- \\
random                 & ---                     & $0.071$ & ---      & --- \\
\bottomrule
\end{tabular}
\end{table}

The policy adds genuine driver signal above the reward filter. On the elite readout it
reaches $0.582$, a $+0.20$ gain over the random-plus-elite null (paired Wilcoxon
$p{=}1.7\times10^{-5}$, $n{=}41$ held-out fates), and on the LOO readout it reaches $0.712$
($+0.29$ above the same null). Beam search reaches $0.507$ ($+0.12$ above the filter,
$p{=}0.02$), so the policy \emph{matches} beam search on absolute recovery. The direct paired
comparison of the policy against beam search is non-significant (parity), consistent with the
matched, paired analysis. The advantage is compute. The policy uses $\approx256$ world-model
evaluations per fate (the LOO readout $\approx207$), while beam search uses $\approx9240$, so
the policy attains the same recovery at roughly $36$--$45\times$ lower inference cost. Both
learned and search methods sit far above the additive ($0.012$) and random ($0.071$)
baselines, which have no world model.

Table~\ref{tab:elite_permetric} extends the comparison across four metrics on the same elite readout and the same
$41$ held-out fates. Beam search (a reward-guided selector read out on its own proposals) and the policy are at
parity on the broad metrics, recall@$20$ and soft-recall@$20$ (paired Wilcoxon $p{=}0.10$ and $0.20$, not
significant), while the policy is significantly ahead at the top of the ranking, recall@$5$ and precision@$5$
(both $p<0.05$); every metric is far above the additive and random baselines ($p<10^{-4}$). Soft-recall credits
functional near-misses through the drift-removed single-TF program space and is scored identically for every
method.

\begin{table}[h]\centering\footnotesize
\caption{Per-metric budget-matched comparison on the elite (top-$10\%$-by-reward) readout, $n{=}41$ held-out
Briggs fates. Baselines are read out on their own reward-selected proposals and restricted to the identical
per-fate shortlist; $p$ is a paired per-fate Wilcoxon of the policy against each baseline.}
\label{tab:elite_permetric}
\begin{tabular}{@{}lcccc@{}}
\toprule
metric & policy & MPC (beam) & additive & random \\
\midrule
recall@$5$        & $0.357$ & $0.235$ ($p{=}0.02$)        & $0.000$ ($p{<}10^{-4}$) & $0.049$ ($p{<}10^{-4}$) \\
recall@$20$       & $0.578$ & $0.507$ (ns, $p{=}0.10$)    & $0.012$ ($p{<}10^{-4}$) & $0.071$ ($p{<}10^{-4}$) \\
soft-recall@$20$  & $0.799$ & $0.716$ (ns, $p{=}0.20$)    & $0.089$ ($p{<}10^{-4}$) & $0.540$ ($p{=}2{\times}10^{-4}$) \\
precision@$5$     & $0.198$ & $0.130$ ($p{=}0.01$)        & $0.000$ ($p{<}10^{-4}$) & $0.020$ ($p{<}10^{-4}$) \\
\bottomrule
\end{tabular}
\end{table}

\subsection{Held-out fate breakdown}\label{app:fatebreakdown}
The Briggs fate universe is $65$ cell-type fates partitioned by a leakage-clean by-name five-fold split (folds
$0$--$4$, $\approx13$ fates each) with a difficulty tier (A $=$ positive-silhouette separable target, B $=$
separable and reachable). The flagship holds out folds $0$, $1$, and $2$, giving the $n{=}41$ held-out fates in
Table~\ref{tab:fatelist}. Fold $0$ was evaluated at a single seed and folds $1$--$2$ at five seeds, so fold-$0$
per-fate values are noisier. The table lists each fate's fold, tier, driver count, and elite recall@$20$ for the
policy (pol) and beam search (MPC), showing the per-fate spread behind the paired test: the policy wins on many
fates, ties on some, and loses on others, which is why the aggregate is a statistical tie with beam search.

\begin{table}[h]\centering\scriptsize
\caption{Per-fate elite recall@$20$ on the $n{=}41$ held-out Briggs fates (folds $0$--$2$).}
\label{tab:fatelist}
\begin{minipage}[t]{0.49\linewidth}\centering
\begin{tabular}{@{}llp{2.6cm}ccc@{}}
\toprule
fold & tier & fate & drv & pol & MPC \\
\midrule
-1 & B & Liver-Hematopoietic stem cells & 4 & 0.25 & 0.50 \\
0 & B & Heart-Epicardial fat cells & 3 & 1.00 & 1.00 \\
0 & A & Spleen-Mesothelial cells & 3 & 1.00 & 1.00 \\
0 & B & Kidney-Ureteric bud cells & 4 & 0.75 & 0.75 \\
0 & B & Cerebellum-Granule neurons & 3 & 0.67 & 0.67 \\
0 & B & Ciliated epithelial cells & 3 & 0.67 & 0.67 \\
0 & A & Eye-PDE11A\_FAM19A2 positive cells & 3 & 0.67 & 0.33 \\
0 & A & Placenta-Trophoblast giant cells & 3 & 0.67 & 0.67 \\
0 & B & Skeletal muscle cells & 3 & 0.67 & 0.67 \\
0 & B & Adrenal-Sympathoblasts & 5 & 0.60 & 0.20 \\
0 & B & Placenta-Syncytiotrophoblasts and villous cytotrophoblasts & 2 & 0.50 & 0.00 \\
0 & B & Placenta-Vascular endothelial cells & 3 & 0.33 & 1.00 \\
0 & B & Vascular endothelial cells & 3 & 0.33 & 0.67 \\
0 & B & Eye-Retinal progenitors and Muller glia & 4 & 0.25 & 0.25 \\
0 & B & Intestine-Stromal cells & 2 & 0.00 & 0.00 \\
1 & B & Pancreas-CCL19\_CCL21 positive cells & 3 & 1.00 & 0.67 \\
1 & B & Placenta-Stromal cells & 2 & 1.00 & 0.50 \\
1 & A & Stomach-Squamous epithelial cells & 3 & 0.93 & 0.67 \\
1 & B & Purkinje neurons & 2 & 0.80 & 0.50 \\
1 & B & Stomach-Parietal and chief cells & 2 & 0.80 & 0.00 \\
1 & B & Schwann cells & 3 & 0.73 & 1.00 \\
\bottomrule
\end{tabular}
\end{minipage}\hfill
\begin{minipage}[t]{0.49\linewidth}\centering
\begin{tabular}{@{}llp{2.6cm}ccc@{}}
\toprule
fold & tier & fate & drv & pol & MPC \\
\midrule
1 & B & Intestine-Smooth muscle cells & 4 & 0.70 & 0.50 \\
1 & B & Eye-Bipolar cells & 3 & 0.67 & 0.87 \\
1 & B & Kidney-Mesangial cells & 3 & 0.60 & 0.00 \\
1 & B & Pancreas-Ductal cells & 3 & 0.60 & 0.00 \\
1 & B & AFP\_ALB positive cells & 3 & 0.53 & 0.33 \\
1 & B & Adrenal-Chromaffin cells & 3 & 0.33 & 0.33 \\
1 & B & Eye-Amacrine cells & 2 & 0.30 & 0.00 \\
2 & A & Heart-Visceral neurons & 1 & 1.00 & 1.00 \\
2 & B & Placenta-IGFBP1\_DKK1 positive cells & 4 & 0.90 & 1.00 \\
2 & B & Lung-Bronchiolar and alveolar epithelial cells & 3 & 0.80 & 0.33 \\
2 & B & Intestine-Intestinal epithelial cells & 2 & 0.50 & 0.50 \\
2 & A & Lung-Vascular endothelial cells & 2 & 0.50 & 0.70 \\
2 & B & Stomach-Goblet cells & 2 & 0.50 & 0.50 \\
2 & B & Pancreas-Islet endocrine cells & 3 & 0.47 & 0.33 \\
2 & B & Eye-Astrocytes & 2 & 0.40 & 0.50 \\
2 & B & Liver-Vascular endothelial cells & 2 & 0.40 & 0.50 \\
2 & A & Chromaffin cells & 4 & 0.35 & 0.75 \\
2 & B & Spleen-Vascular endothelial cells & 3 & 0.33 & 0.67 \\
2 & B & Cerebellum-Inhibitory interneurons & 4 & 0.20 & 0.25 \\
2 & B & Adrenal-Vascular endothelial cells & 1 & 0.00 & 0.00 \\
\bottomrule
\end{tabular}
\end{minipage}
\end{table}

\subsection{Axis 2, diversity at quality}
The second axis asks whether the planner returns \emph{many} high-reward cocktails rather than
one. On the enumerable HSPC candidate pool we draw $256$ proposals from each method and count
how many are distinct and above a fixed quality bar (Table~\ref{tab:diversity}). Asked to
return $256$ candidates, budget-matched beam search collapses to a single distinct high-quality
cocktail. Size-matched random sampling gives $64$ distinct sets but of low reward. The policy
returns $154$ distinct high-quality cocktails. The policy is the only method that is both
diverse and high-quality. This diversity is owned by the goal-conditioned $\log Z$. It lets
the effective number of distinct sampled sets approach the target ceiling ($\mathrm{ess}_q$ up
to $\approx46$), whereas the shared-$\log Z$ variant collapses to an amortized argmax
($\mathrm{ess}_q\approx1$--$2$).

\begin{table}[h]\centering\footnotesize
\caption{Diversity at quality on the enumerable HSPC pool, $256$ proposals per method.
Distinct high-quality count is the number of distinct cocktails above a fixed reward bar.
Beam search is asked for $256$ candidates and returns its frontier. It is not a sampler.}
\label{tab:diversity}
\begin{tabular}{@{}lcc@{}}
\toprule
Method & distinct high-quality cocktails & quality \\
\midrule
policy (flagship)      & $154$ & high \\
beam-top-$256$         & $1$   & high \\
random (size-matched)  & $64$  & low  \\
\bottomrule
\end{tabular}
\end{table}

\subsection{Zero-shot HSPC, coverage advantage and honest limits}
We test the frozen policy zero-shot on a held-out hESC-to-HSPC fate. This is a single fate, so we
report it as a case study over the $11$ fold$\times$seed runs with no per-fate significance test.

\textbf{Diversity-coverage advantage on the arrayed key.} Against the Briggs arrayed answer key
($54$ driver TFs across $28$ arrayed-validated combinations from the supplementary screen),
the policy's diverse elite proposals cover far more of the key than the search or heuristic
baselines, an elite recall@$50$ of $0.131$ against $\le0.056$ for MPC and additive and $0.044$ for
random, and $15$ of the $54$ drivers surfaced across the runs against $3$--$4$ for every baseline
(Table~\ref{tab:hspc}). The surfaced TFs are biologically coherent: the policy proposes the
canonical hematopoietic regulators LYL1, KLF1, KLF2, and IKZF1 in $8$--$11$ of the $11$ runs, and
members of the blood stem-cell heptad (TAL1/SCL, GATA2, RUNX1, FLI1) and the self-renewal factor
HOXB4 in $1$--$3$ runs, all with the regulatory prior switched off.

\textbf{Honest limits.} The advantage is coverage from diverse best-of-$N$ sampling, not per-driver
precision. On the narrow twelve-TF $P_1{\cup}P_2$ named key, per-run elite recall@$20$ is $0.12$ for
the policy against $0.27$ for MPC and $0.15$ for its random-plus-elite null, so the policy does not
beat search or its own reward-filtered null on that key. We do not calibrate a breadth-matched null
for the arrayed key, so the arrayed advantage is established against the search baselines but not
against equal-breadth random sampling. We therefore treat HSPC as a diversity-coverage result on the
broad key, not a beyond-chance per-driver recovery, and note that the reward filter, rather than the
policy alone, contributes to surfacing canonical TFs. Zero-shot transfer to a new fate is a
data-coverage limit, addressable with in-domain fate data.

\begin{table}[h]\centering\footnotesize
\caption{Zero-shot HSPC coverage of the $54$-TF Briggs arrayed answer key, elite readout, over $11$
fold$\times$seed runs (single fate, no p-value). Baselines are read out on their own reward-selected
proposals. ``surfaced'' is the number of distinct arrayed driver TFs appearing across the runs.}
\label{tab:hspc}
\begin{tabular}{@{}lcc@{}}
\toprule
Method & recall@$50$ & distinct drivers surfaced (of $54$) \\
\midrule
policy (flagship)        & $0.131$ & $15$ \\
MPC / greedy             & $0.056$ & $3$ \\
additive                 & $0.056$ & $3$ \\
random shortlist         & $0.044$ & $4$ \\
\bottomrule
\end{tabular}
\end{table}

\subsection{Honest negatives}
Two engineering attempts did not help and we report them. Prioritized, reward-weighted replay
hurt relative to plain recency-weighted replay, so the deployed GFlowNet uses plain
recency-weighted replay. Bulk-signal removal did not lift the cross-fate discrimination wall.
The policy's own-minus-other discrimination sits near zero while a drivers-as-picks oracle
reaches about $0.65$ on Briggs, and the wall persists across reward shaping, exploration
budget, and bulk removal, so we report it as a limit of designing against signature-defined
fates rather than a bug. The SOFT functional-space metrics are saturated as described in
App.~\ref{app:metrics}. On Briggs the policy SOFT-1 is about $0.28$ against a null of about
$0.28$, essentially no separation, which is why SOFT is not a headline.

\subsection{Prediction leakage ladder}
The world-model prediction advantage over the fair additive baseline survives progressively
harder holdouts. We retrained the held-out finetune under two extra rungs, leave-TF-out
(rung 2) and leave-program-out (rung 3, leave-$P_1$-out), and re-measured
drift-removed \PearsonDEG. Table~\ref{tab:ladder} shows that on held-out Briggs
combinations the world model with its interaction head beats both the observed-additive and
the model-additive baselines at every rung, and the margin holds at high order $k6$--$8$.
The honest counterpoint is the answer-key row. The $222$ high-order answer-key combinations
are extrapolations the world model does not predict better than observed-additive
($0.133$ versus $0.208$ at rung2), consistent with the main-text limitation that
composition is effectively additive for held-out HSPC TFs that never appear in the training
combinations.

\begin{table}[h]\centering\footnotesize
\caption{Prediction leakage ladder, drift-removed \PearsonDEG, world model with head
versus observed-additive and model-additive baselines. Rung labels are the
progressively harder retrain holdouts. Values are group means with $n$ combinations.}
\label{tab:ladder}
\begin{tabular}{@{}llcccc@{}}
\toprule
Rung / group & $n$ & WM$+$head & WM no-head & obs-additive & model-additive \\
\midrule
\multicolumn{6}{@{}l}{\textbf{Rung2, leave-TF-out retrain}}\\
Briggs combos train (in-dist)   & $387$ & $0.623$ & $0.385$ & $0.206$ & $0.294$ \\
Briggs combos held-out (set-out)& $97$  & $0.464$ & $0.358$ & $0.245$ & $0.280$ \\
Briggs held-out $k6$--$8$       & $6$   & $0.478$ & $0.087$ & $0.096$ & $0.064$ \\
Joung combos held-out           & $8$   & $0.957$ & $0.397$ & $0.134$ & $0.100$ \\
Briggs answer-key (extrapol.)   & $222$ & $0.133$ & $0.119$ & $0.208$ & $0.111$ \\
\multicolumn{6}{@{}l}{\textbf{Rung3, leave-program-out retrain}}\\
Briggs combos train (in-dist)   & $602$ & $0.591$ & $0.381$ & $0.197$ & $0.259$ \\
Briggs combos held-out          & $151$ & $0.459$ & $0.335$ & $0.224$ & $0.231$ \\
Briggs held-out $k6$--$8$       & $26$  & $0.418$ & $0.115$ & $0.122$ & $0.054$ \\
\bottomrule
\end{tabular}
\end{table}

% =============================================================================
\section{Additional world-model results}\label{app:wm}
\subsection{Reconstruction, discrimination, and per-DEG behavior}
The deployed production model reaches a stable reconstruction MSE of about $0.52$ and a KL
of about $0.26$ at epoch $60$ with the acyclicity residual at zero throughout (training
log). Drift-removed discrimination on Joung singles is near ceiling in the finetune
checkpoint, with held-out singles reaching AUROC $0.917$ at rung2 and $0.952$ at rung3 while
the observed-additive control is $1.0$ by construction. The world model's advantage grows
with combination order. On held-out Briggs combinations the model-with-head \PearsonDEG
stays high at $k6$--$8$ ($0.478$ at rung2) while both additive baselines collapse
toward $0.06$ to $0.12$ (Table~\ref{tab:ladder}), which is the signature of genuine
non-additive structure at high order.

\subsection{The directed causal graph is load-bearing for prediction}\label{app:causal_load}
The causal-structure panels of Fig.~\ref{fig:epi} (E, F) in the main text are the result this
section expands. We ask a sharp, testable question about the learned graph $A$, namely whether its
\emph{direction} is load-bearing for prediction. We answer by ablating $A$ inside the frozen do-operator
(reversing, symmetrizing, or shuffling its edges) and re-measuring $\Delta$-prediction on the
deployed production world model. This is a relative true-vs-shuffled comparison on
the same perturbations, so it needs no held-out split.

Panels A and B show single-TF and combination prediction under original, reversed,
symmetrized, and shuffled $A$, scored by drift-removed \PearsonDEG and cross-fate
discrimination AUROC. Against a $1000$-shuffle null, single-TF prediction collapses from
Pearson $0.536\,[0.443,0.625]$ to a shuffled mean of $-0.156\pm0.010$, a gap of $+0.692$ that
no shuffle reaches (permutation $p<0.001$, paired Wilcoxon across $66$ TFs
$p{=}8.9\times10^{-11}$, with $59$ of $66$ TFs predicting worse under a shuffled graph). Discrimination
AUROC falls from $0.84$ to $0.60$ and top-1 driver recall from $0.73$ to $0.02$. Reversing the
edges is as damaging as shuffling them, so the effect is directional rather than a matter of
connectivity alone. On combinations the gap is smaller ($+0.102$, permutation $p<0.001$,
Wilcoxon $p{=}6.1\times10^{-5}$, all $14$ combinations improving) because the drift-removed
combo response is dominated by a shared high-order shift. We therefore quote the single-TF
effect as the headline. The full ablation numbers are in Table~\ref{tab:edge}.

Panel C shows that the latent graph is biologically structured. The masked decoder pins each
latent axis to a named GO:BP program at mask coherence $1.00$, and the highest-connectivity
hubs of $A$ are developmental programs. We do \emph{not} claim $A$ recovers a curated GRN.
With $A{=}0$ the do-operator injects only the direct shift $S_g$, which already ranks a TF's
direct targets, so $A{=}0$ stays competitive on a direct-target metric. Propagation over the
directed $A$ is what carries the genuine downstream and combination response, which is where
the directed graph is load-bearing.

\subsection{The edge-direction ablation in full}
Table~\ref{tab:edge} gives the complete edge-direction ablation, computed on the deployed
production world model. For single TFs the drift-removed \PearsonDEG collapses
from $0.536$ with $95\%$ bootstrap CI $[0.443,0.625]$ to a shuffled-$A$ mean of
$-0.156\pm0.010$ over $1000$ shuffles, a gap of $+0.692$ with permutation $p<0.001$ (no
shuffle exceeds the original) and paired Wilcoxon $p=8.9\times10^{-11}$ across $66$ TFs.
Reversing the edges ($-0.130$) is about as damaging as shuffling them, which shows the
effect is directional and not merely connectivity. Top-1 driver recall collapses from
$0.727$ to $0.015$ and discrimination AUROC from $0.844$ to $0.600$. On combinations the raw
Pearson is drift-dominated (original $0.977$), so the gap is smaller at $+0.102$ but still
significant (permutation $p<0.001$, Wilcoxon $p=6.1\times10^{-5}$ across $14$ combinations,
all improve). We therefore quote the single-TF gap as the headline and do not carry $+0.69$
over to combinations.

\begin{table}[h]\centering\footnotesize
\caption{Edge-direction ablation on the deployed model, drift-removed \PearsonDEG,
top-1 driver recall, and discrimination AUROC under original, reversed, symmetrized, and
shuffled $A$. Non-original variants are rescaled to match $\lVert A\rVert_2$.}
\label{tab:edge}
\begin{tabular}{@{}lcccc@{}}
\toprule
 & original & reversed & symmetrized & shuffled \\
\midrule
Single Pearson ($n{=}66$)   & $0.536$ & $-0.130$ & $0.107$ & $-0.156$ \\
Single recall@1             & $0.727$ & $0.015$  & $0.076$ & $0.015$ \\
Single disc AUROC           & $0.844$ & $0.592$  & $0.686$ & $0.600$ \\
Combo Pearson ($n{=}14$)    & $0.977$ & $0.797$  & $0.918$ & $0.875$ \\
Combo recall@1              & $0.429$ & $0.071$  & $0.071$ & $0.057$ \\
\bottomrule
\end{tabular}
\end{table}

\subsection{The retired ``$+0.68$'' result}
An earlier combination-prediction figure of about $+0.68$ over additive came from a
non-deployable checkpoint and a raw, non-drift-removed Pearson against a cross-context
additive baseline that itself scored near zero. That comparison conflated cross-context
drift with model skill, because a large part of the raw combination Pearson is the shared
mean shift that any model captures. Once we remove drift and compare against a fair,
same-context additive baseline, the honest advantage is the $+0.18$ to $+0.23$ reported in
the main text and corroborated by the leakage ladder in Table~\ref{tab:ladder}. We retire
the $+0.68$ figure.

\subsection{Magnitude over-prediction and calibration}
The interaction sign and ranking are recovered but the interaction magnitude is
uncalibrated. The DEG-MSE term is deliberately unwhitened so the model does track per-gene
response magnitude, yet on epistasis the recovered interaction carries a sub-additive
offset rather than a clean scale, so we report sign and rank recovery and flag the magnitude
as an honest calibration limit. This matches the leave-one-out cross-validated pairwise
interaction correlation of $+0.39\pm0.088$ reported in the main text, which measures rank
and sign agreement rather than calibrated magnitude.

\subsection{Identifiability and rotational ambiguity}
Fixing the decoder to GO programs pins each latent axis to a named biological process and
removes the coordinate permutation ambiguity of the latent, following the identifiable
causal-disentanglement lineage of Discrepancy-VAE and its GO-structured variant
\citep{zhang2023discrepancyvae,sena2025}. We do not claim full structural identifiability of
$A$. The effective decoder that maps the exogenous term to genes is
$D_{\mathrm{eff}}=D\big(I-\operatorname{diag}(1-u)A^{\top}\big)^{-1}$, and any invertible
transform $T$ with $A'^{\top}=T A^{\top}T^{-1}$ and $D'=DT^{-1}$ that preserves the masked
structure leaves the observational and interventional distributions unchanged. This is a
rotational ambiguity in $A$, which we treat as a known caveat. Our load-bearing claim is the
weaker and testable one. The graph direction improves prediction under the edge-direction
ablation (Table~\ref{tab:edge}), and the do-operator is exact for composition
(Eq.~\ref{eq:doop}). We do not claim that $A$ is a recovered gene-regulatory network. Because
the model is trained prior-off on only TFs, TF-to-target databases are off-predictors in our
regime, and a plain co-expression baseline edges out the projected graph on every reference,
as expected for a context-specific prior-off model \citep{genernib2025,collectri2023,garcia2019dorothea,szklarczyk2023string}.
